\documentclass[
    a0paper,
    landscape,
    margin=0mm,
    innermargin=14mm,
    blockverticalspace=2mm,
    colspace=10mm
]{tikzposter}

\usepackage[T1]{fontenc}
\usepackage{amsmath}
\usepackage{amssymb}
\usepackage{booktabs}
\usepackage{array}
\usepackage{graphicx}
\usepackage{xcolor}
\usepackage{ragged2e}
\usepackage{hyperref}
\usepackage{url}
\usepackage{enumitem}

\renewcommand{\familydefault}{\rmdefault}
\setlength{\parindent}{0pt}
\setlength{\parskip}{0.15cm}
\linespread{1.00}

\setlist[itemize]{
    leftmargin=1.3em,
    itemsep=0.10em,
    topsep=0.15em,
    parsep=0pt
}

% ---------- Colour palette ----------
\definecolor{PosterGreen}{RGB}{45,105,55}
\definecolor{PosterText}{RGB}{35,35,35}
\definecolor{PosterGray}{RGB}{100,100,100}

\usetheme{Simple}
\tikzposterlatexaffectionproofoff
\titlegraphic{}
\useblockstyle{Minimal}

\definecolor{PosterBackground}{RGB}{232,242,233}
\colorlet{backgroundcolor}{PosterBackground}
\colorlet{titlebgcolor}{white}
\colorlet{titlefgcolor}{PosterText}
\colorlet{blocktitlebgcolor}{white}
\colorlet{blocktitlefgcolor}{PosterGreen}
\colorlet{blockbodybgcolor}{white}
\colorlet{blockbodyfgcolor}{PosterText}
\colorlet{framecolor}{white}

\definetitlestyle{CleanAcademicTitle}{
    width=\paperwidth,
    roundedcorners=0,
    linewidth=0pt,
    innersep=8mm,
    titletotopverticalspace=5mm,
    titletoblockverticalspace=5mm
}{}
\usetitlestyle{CleanAcademicTitle}

\renewcommand{\arraystretch}{1.0}

\hypersetup{
    colorlinks=true,
    urlcolor=PosterGreen
}

\newcommand{\headrule}{%
    \par\vspace{-0.55cm}%
    {\color{PosterGreen}\rule{0.55\linewidth}{0.7pt}}%
    \par\vspace{0.18cm}%
}
\newcommand{\captiontext}[1]{%
    \par\vspace{0.12cm}%
    {\centering\normalfont\color{PosterGray}#1\par}%
    \vspace{0.08cm}%
}

\newcommand{\screenshot}[2]{%
    \setlength{\fboxsep}{2pt}%
    \fcolorbox{PosterGreen}{white}{%
        \includegraphics[
            width=0.92\linewidth,
            height=#2,
            keepaspectratio
        ]{#1}%
    }%
}


% ---------- Title ----------
\title{
    \fontsize{52pt}{56pt}\selectfont
    From-Scratch \(\arccos(x)\) Calculator:
    Taylor Series, Newton's Method, and Software Quality Assurance
}

\author{
    \Large\bfseries Mohammad Aliyawar Khan \quad|\quad Student ID: 40309082
}

\institute{
    \normalsize
    Department of Computer Science and Software Engineering,
    Concordia University\\[0.05cm]
    SOEN 6011 -- Deliverable 3 \quad|\quad Summer 2026
}

\begin{document}
\maketitle

\begin{columns}

% =========================================================
% COLUMN 1 -- BACKGROUND AND NUMERICAL METHOD
% =========================================================
\column{0.245}

\block{D2 Baseline: Numerical Method}{
\headrule
\justifying
The calculator computes the principal real value of $\arccos(x)$ for
$-1\leq x\leq1$, with results in $[0,\pi]$. It uses range reduction,
a Taylor-series recurrence for $\arcsin$, and Newton's method for
square roots without production use of \texttt{math}.

\[
\arccos(x)=
\begin{cases}
2\arcsin\left(\sqrt{(1-x)/2}\right), & x\geq0,\\
\pi-\arccos(-x), & x<0.
\end{cases}
\]

Exact endpoint handling returns $\arccos(1)=0$ and
$\arccos(-1)=\pi$. The tested error target is
$\leq10^{-6}$ radians.
}

\block{D3 Problem 7: Software Quality and Accessibility}{
\headrule
\justifying
Problem 7 improves the D2 calculator through PEP~8, Flake8, Pylint,
\texttt{pdb}, Semantic Versioning, UIDP-guided interface design, and
accessibility-oriented improvements.

The production implementation remains from scratch: it does not import
\texttt{math} and does not use library inverse-trigonometric or
square-root functions.
}

\block{GUI Demonstration}{
\headrule
\begin{center}
\screenshot{gui-success.png}{6.0cm}
\captiontext{
Exact endpoint calculation:
$\arccos(-1)=\pi=3.1415926536$ radians.
}
\end{center}
}
\vspace{-0.35cm}

\block{Implementation and Versioning}{
\headrule
\justifying
\textbf{From-scratch implementation.} The application manually
implements absolute value, square root, arcsine, and arccosine.

Tkinter is used only for the interface; \texttt{math} is not imported
by \texttt{problem7.py}.

\textbf{Validation and exceptions}
\begin{itemize}
    \item \texttt{InputValidationError}: NaN, infinity, or invalid domain
    \item \texttt{ConvergenceError}: configured iteration limit reached
\end{itemize}

\textbf{Semantic Versioning}
\[
\boxed{\texttt{v1.0.0 (D2)} \longrightarrow
       \texttt{v1.1.0 (D3)}}
\]

\textbf{Why minor?} D3 adds backward-compatible quality and usability
features: automated tests, static analysis, debugger evidence,
keyboard support, visible focus, and clearer status feedback. The
public calculator behavior and valid input domain are unchanged.

The identifier \texttt{v1.1.0} is consistent in the code, GUI title,
README, and Git tag.
}
% =========================================================
% COLUMN 2 -- SOFTWARE QUALITY TOOLS
% =========================================================
\column{0.245}

\block{PEP 8 and Flake8}{
\headrule
\justifying
Applied practices include four-space indentation, descriptive names,
docstrings, grouped constants, wrapped expressions, and separation of
numerical and GUI responsibilities.

\begin{center}
\ttfamily
python -m flake8 problem7.py test\_problem7.py
\end{center}

\textbf{Analyzed files:} \texttt{problem7.py} and
\texttt{test\_problem7.py}.\\
\textbf{Final result:} \textit{0 violations reported.}

\begin{center}
\screenshot{flake8-result.png}{4.5cm}
\captiontext{Flake8 execution evidence for the final source and test files.}
\end{center}
}

\block{Static Analysis with Pylint}{
\headrule
\justifying
Pylint inspected the final source and test files for errors, naming,
documentation, unused code, exception handling, and maintainability.

\begin{center}
\ttfamily
python -m pylint problem7.py test\_problem7.py
\end{center}

\textbf{Analyzed files:} \texttt{problem7.py} and
\texttt{test\_problem7.py}.\\
\textbf{Final score:} \texttt{10.00/10}.\\
\textbf{Remaining messages:} None.

\begin{center}
\screenshot{pylint-result.png}{4.5cm}
\captiontext{Pylint report for the final source and test files.}
\end{center}
}


\block{Debugging with \texttt{pdb}}{
\headrule
\justifying
A breakpoint was placed in \texttt{calculate\_arccos()} immediately
before range reduction for $x=0.9999$. The session used \texttt{n},
\texttt{s}, \texttt{p}, and \texttt{c} to inspect \texttt{x},
\texttt{reduced\_input}, \texttt{term}, \texttt{coefficient}, and
\texttt{result}.

\textbf{Confirmed:} range reduction produces a small arcsine argument
near the endpoint; Taylor terms decrease to the configured tolerance;
and the final value remains in the principal range $[0,\pi]$. The
investigation found no numerical defect.

\begin{center}
\begin{tabular}{ll}
\toprule
\textbf{Command} & \textbf{Purpose}\\
\midrule
\texttt{n} & Execute the next line\\
\texttt{s} & Step into a function\\
\texttt{p variable} & Inspect a variable\\
\texttt{c} & Continue execution\\
\bottomrule
\end{tabular}
\end{center}

\begin{center}
\screenshot{pdb-session-a.png}{3.5cm}
\captiontext{Breakpoint, validation, and range-reduction inspection.}
\vspace{0.05cm}
\screenshot{pdb-session-b.png}{3.5cm}
\captiontext{Taylor recurrence variables and final-result inspection.}
\end{center}
}

% =========================================================
% COLUMN 3 -- DEBUGGING AND TESTING
% =========================================================
\column{0.265}

\block{D3 Problem 8: Unit Testing}{
\headrule
\justifying
The automated suite uses Python's built-in \texttt{unittest} framework
(PyUnit) and imports the D3 production module \texttt{problem7.py}.

\textbf{Coverage:} helper functions, Newton square root, Taylor arcsine,
exact endpoints, normal inputs, near-endpoint inputs, invalid-domain
values, NaN, infinity, principal range, numerical accuracy, and
exception paths.

\begin{center}
\ttfamily
python -m unittest -v test\_problem7.py
\end{center}

\textbf{Result: 31 tests passed.} Reference values from
\texttt{math.acos()} and \texttt{math.asin()} are test oracles only;
the production module \texttt{problem7.py} does not import
\texttt{math}.

\begin{center}
\screenshot{testcase1.png}{3.5cm}
\captiontext{PyUnit evidence: first group of test cases.}
\vspace{0.06cm}
\screenshot{testcase2.png}{3.5cm}
\captiontext{PyUnit evidence: final tests and \texttt{OK} result.}
\end{center}
}
\block{Selected Test Cases}{
\headrule
\justifying
\begin{center}
\begin{tabular}{p{0.12\linewidth}p{0.25\linewidth}p{0.50\linewidth}}
\toprule
\textbf{ID} & \textbf{Input} & \textbf{Expected outcome}\\
\midrule
TC01 & $1.0$ & Exact result $0.0$\\
TC02 & $-1.0$ & Exact result $\pi$\\
TC03 & $0.0$ & Approximately $\pi/2$\\
TC04 & $0.5$ and $-0.5$ & Approximately $\pi/3$ and $2\pi/3$\\
TC05 & $\pm0.9999$ & Accurate near-endpoint result\\
TC06 & $1.1$, $-1.1$ & \texttt{InputValidationError}\\
TC07 & NaN, $\pm\infty$ & \texttt{InputValidationError}\\
TC08 & Valid test set & Error $\leq10^{-6}$ radians\\
\bottomrule
\end{tabular}
\end{center}
}
\block{UIDP and Accessibility}{
\headrule
\small
\justifying

\textbf{Directly applied}
\begin{itemize}
    \item \textbf{User Profiling:} supports users entering one real
    value $x$ and expecting an angle in radians.
    \item \textbf{Humility:} clear guidance and recoverable errors avoid
    blaming the user.
    \item \textbf{Feature Exposure:} labelled input plus visible
    Calculate, Clear, and Exit controls.
    \item \textbf{State Visualization:} separate result and status areas
    show ready, success, and error states.
    \item \textbf{Coherence:} consistent labels, colours, shortcuts, and
    error wording.
    \item \textbf{Safety:} validation prevents invalid calculations while
    preserving input for correction.
\end{itemize}

\textbf{Limited applicability:} \textbf{Metaphor} is limited because a scientific calculator is already a familiar direct-manipulation model.

\textbf{Accessibility aim:} readable text, high contrast, visible focus, keyboard navigation, shortcuts, and textual feedback not based only on
colour. The GUI \emph{aims to be accessible}; it was not formally evaluated with users or assistive technology.

\begin{center}
\screenshot{uidp-mind-map.png}{6.0cm}
\captiontext{UIDP applicability decision mind map.}
\end{center}
}

% =========================================================
% COLUMN 4 -- GAI, CONCLUSION, REFERENCES
% =========================================================
\column{0.245}
\block{GAI and CASTROFF Evidence}{
\headrule
\justifying

\textbf{AI tool:} Perplexity AI.

\textbf{Problem 7: Code Quality and Accessibility}\\
\textbf{Prompt:} ``Review my from-scratch Python Tkinter
$\arccos(x)$ calculator for PEP~8, maintainability, UIDP, and
accessibility. Do not recommend prohibited \texttt{math} functions.''

\textbf{CASTROFF:} \textit{Constraints}---from-scratch code, no
production \texttt{math}; \textit{Audience}---SOEN 6011 evaluator;
\textit{Structure}---prioritized review; \textit{Tone}---technical;
\textit{Role}---quality reviewer; \textit{Output}---actionable
recommendations; \textit{Focus}---PEP~8, UIDP, and accessibility;
\textit{Function}---improve D3 quality.

\textbf{Outcome:} Accepted clearer feedback, visible focus, keyboard
support, and documentation. Rejected suggestions violating the
from-scratch constraint. Recommendations were checked against the code
and compiled GUI.

\textbf{Problem 8: Boundary and Exception Testing}\\
\textbf{Prompt:} ``Review my PyUnit tests and suggest boundary,
invalid-input, exception, convergence, and accuracy tests. Use
\texttt{math.acos()} and \texttt{math.asin()} only as test oracles.''

\textbf{CASTROFF:} \textit{Constraints}---existing API, no production
\texttt{math}; \textit{Audience}---SOEN 6011 evaluator;
\textit{Structure}---boundary and exception tests; \textit{Tone}---technical;
\textit{Role}---test reviewer; \textit{Output}---executable PyUnit tests;
\textit{Focus}---boundaries, exceptions, accuracy, and range;
\textit{Function}---validate the implementation.

\textbf{Outcome:} Tests were reviewed against the implementation and
executed locally: \textbf{31 tests passed}. The \texttt{math} functions
were used only as test oracles, not production dependencies.
}

\block{Conclusion}{
\headrule
\justifying
The D3 release combines a range-reduced Taylor series and Newton square
root with validation, exception recovery, and a Tkinter GUI. Flake8,
Pylint, \texttt{pdb}, Semantic Versioning, UIDP-guided design, and
automated PyUnit tests provide verifiable quality evidence.

\textbf{Result:} a numerically reliable, maintainable, testable, and
accessibility-oriented $\arccos(x)$ calculator.
}

\block{References and Repository}{
\headrule
\small
\justifying

\begin{tabular}{@{}r@{\hspace{0.25em}}p{0.84\linewidth}@{}}
[1] & Concordia University, \emph{SOEN 6011 Project Description},
Summer 2026. \\

[2] & Python Software Foundation, \emph{PEP 8 -- Style Guide for
Python Code}. \\

[3] & PyCQA, \emph{Flake8 Documentation}. \\

[4] & Pylint Contributors, \emph{Pylint Documentation}. \\

[5] & Python Software Foundation, \emph{pdb -- The Python Debugger}. \\

[6] & Python Software Foundation, \emph{unittest -- Unit Testing Framework}. \\

[7] & Semantic Versioning, \emph{Semantic Versioning 2.0.0}. \\

[8] & W3C, \emph{Web Content Accessibility Guidelines (WCAG) 2.2}. \\

[9] & SOEN 6011 Course Material, \emph{CASTROFF Prompt Framework}. \\
\end{tabular}

\vspace{0.10cm}

\normalsize\textbf{Repository:}\\[0.02cm]
\ttfamily\small
github.com/md-aliyawar-khan/SOEN-6011-D3-arccos-calculator
}

\end{columns}
\end{document}